\subsection{Learned stopping rule}
Every 5\,ms, the stopping rule compares the predicted gain of searching for another $H=10$\,ms with the cost of doing so.
With $\widehat{\Delta J}$ as the predicted gain, the search stops at time $\tau$ if
\begin{equation}\label{eq:Sstop}
\widehat{\Delta J}\;\le\;C_{\mathrm{delay}}(\tau)+\mu W H .
\end{equation}
The gain is predicted from the search history by a Poisson generalised linear model with a log
link~\cite{mccullagh1989glm} and a small $L_2$ penalty. Its eleven features describe the elapsed time, the number of
iterations, the recent improvements over 5\,ms, over 20\,ms and over the second half of the elapsed time, the time since
the last improvement, the total improvement over the repaired plan, the unassigned share of the reward, the numbers of
tasks and UAVs, and the iteration rate. The target is the improvement in the next 10\,ms in percent of $W$. The model is
fitted on 500\,ms search trajectories of development scenarios only, and a prediction costs one dot product. The search
loop enforces the hard time budget $B$ independently of the stopping rule. The price of computation is
$\mu=10^{-5}$ per millisecond.

\subsection{Deterministic delay cost}
The delay cost is computed from the best plan found so far. A late start of a route is partly absorbed by the waiting
times along it, in the same way as the forward time slack of routing with time windows~\cite{savelsbergh1992vrptw}. The
tolerance $\mathrm{tol}_j$ of a task is the margin between its finish time and its deadline, plus all waiting times on
the route up to this task. A task is lost when the start shift $\sigma_j(\tau)$ of its route exceeds its tolerance. The
start shift is zero for a UAV that is still busy, and it equals the waiting time until the publication for an idle UAV.
With $v_j=\alpha w_j-\delta\,\mathbb{1}[\text{$j$ changes UAV}]$ as the value of task $j$ in the objective and
$\Lambda(\tau)=\sum_j v_j\,\mathbb{1}[\sigma_j(\tau)>\mathrm{tol}_j]$, the delay cost of searching $H$ milliseconds longer
is $C_{\mathrm{delay}}(\tau)=\Lambda(\tau+H)-\Lambda(\tau)$.

\begin{prop}\label{prop:delay}
For a feasible plan executed exactly, (i) $\Lambda(\tau)$ is the value lost through missed deadlines if every UAV
follows its route, (ii) the loss of the objective is at most $\Lambda(\tau)$ if late tasks are skipped, and
(iii) $\Lambda(\tau)$ does not decrease with $\tau$ and is zero while all UAVs are busy with their committed tasks.
Hence $C_{\mathrm{delay}}$ is non-negative and vanishes while no UAV is waiting.
\end{prop}
\begin{proof}
A route that starts later by some shift reaches its first task later by the same amount. Each waiting time absorbs part
of the delay, so the delay at the end of a task equals the start shift minus all earlier waiting times, if this
difference is positive. The task therefore misses its deadline exactly when the start shift exceeds its tolerance.
Waiting consumes no energy, which proves (i). If late tasks are skipped, the triangle inequality ensures that every
later task is reached no later and with no more energy. All tasks that are on time in (i) remain on time, and the flown
distance, hence the energy term $\gamma E$, does not grow. The objective therefore loses the values $v_j$ of the skipped
tasks and nothing else, which is at most $\Lambda(\tau)$ and proves (ii). Part (iii) holds because the start shift does
not decrease with $\tau$ and every tolerance is non-negative in a feasible plan.
\end{proof}

\subsection{Execution-aware delay cost}\label{sec:Srisk}
Under uncertain execution, task $j$ is lost when $\sigma_j(\tau)+e_j>\mathrm{tol}_j$, where $e_j$ is its finish-time
error. The execution-aware cost replaces the indicator in $\Lambda$ by the probability of this event. The error of task
$j$ is modelled as $e_j\sim\mathcal{N}(m(h_j),s(h_j)^2)$, where $h_j=f_j-t_0$ is its lead time, the nominal finish time
minus the time of the planning snapshot. The mean and the spread are linear in the lead time,
\begin{equation}\label{eq:Serrmodel}
m(h)=b_0+b_1h,\qquad s(h)=s_0+s_1h,
\end{equation}
because errors accumulate along a route. The expected dissemination latency $\bar L$ is added to the start shift. The
expected loss and the execution-aware delay cost are
\begin{equation}\label{eq:Srisk}
\Lambda_{\mathrm{R}}(\tau)=\sum_j v_j\,\Phi\!\left(\frac{\sigma_j(\tau)+\bar L+m(h_j)-\mathrm{tol}_j}{s(h_j)}\right),
\qquad C_{\mathrm{R}}(\tau)=\Lambda_{\mathrm{R}}(\tau+H)-\Lambda_{\mathrm{R}}(\tau),
\end{equation}
where $\Phi$ is the standard normal distribution function. Search and predictor are unchanged, and $C_{\mathrm{R}}$
replaces $C_{\mathrm{delay}}$ in Eq.~\eqref{eq:Sstop}.

\begin{prop}\label{prop:risk}
(i) With $b_0=b_1=s_0=s_1=0$ and $\bar L=0$, $\Lambda_{\mathrm{R}}=\Lambda$, and the stopping rule takes exactly the
decisions of the deterministic one. (ii) $\Lambda_{\mathrm{R}}(\tau)$ does not decrease with $\tau$, so $C_{\mathrm{R}}\ge0$.
(iii) Under the error model, $\Lambda_{\mathrm{R}}(\tau)$ is the expected value of the tasks whose deadline is missed if
every UAV follows its route.
\end{prop}
\begin{proof}
With zero spread, $\Phi$ is read as the indicator of a positive argument, which gives (i). The start shift does not
decrease with $\tau$ and $\Phi$ is non-decreasing, which gives (ii). Part (iii) follows from the proof of
Proposition~\ref{prop:delay}, because task $j$ misses its deadline exactly when its delay at the finish,
$\sigma_j+e_j$ minus the absorbed waiting time, is positive, and the expectation of the indicator of this event is the
probability in Eq.~\eqref{eq:Srisk}.
\end{proof}

The parameters of Eq.~\eqref{eq:Serrmodel} are fitted on development missions. The mean is estimated by least squares,
and the spread by least squares on the absolute residuals scaled by $\sqrt{\pi/2}$, which is less sensitive to the heavy
tail of the errors than squared residuals. The latency $\bar L$ is the mean of the latency distribution in tier B and
the measured mean plan-acceptance latency in tier C.

